\documentclass[11pt]{article}
\usepackage[utf8]{inputenc}
\usepackage[T1]{fontenc}
\usepackage{amsmath,amssymb,amsthm}
\usepackage{hyperref}
\usepackage{listings}

\title{Categorical Library in Anders}
\author{Groupoid Infinity}
\date{May 2026}

\newcommand{\catC}{\mathcal{C}}
\newcommand{\catD}{\mathcal{D}}
\newcommand{\op}{^\text{op}}

\begin{document}

\maketitle

\begin{abstract}
This article describes the design and implementation of the categorical library in the Anders proof assistant.
We focus on incorporating the design principles from Jacques Carette et al. (2021) formalization in Agda, specifically regarding definitional duality and the use of cubical foundations for univalent categories.
\end{abstract}

\section{Introduction}

The Anders categorical library is designed to be a minimal yet powerful implementation of category theory
in a pure Martin-Löf Type Theory (MLTT) with cubical extensions. The library prioritizes
proof-relevance, univalence, and duality-friendliness.

\section{Definitional Duality}

A key design choice, following \cite{carette2021}, is the use of symmetric laws for associativity and identity.
By including both \texttt{assoc} and \texttt{assoc-inv}, as well as \texttt{id-id} (the neutral identity law),
we ensure that the opposite category operation $(\cdot)\op$ is a definitional involution:
\begin{equation}
    (\catC\op)\op \equiv \catC
\end{equation}
This allows for "duality for free," where any theorem proved for a category $\catC$ can be immediately applied
to its dual without explicit transport along an isomorphism.

\section{Cubical Foundations}

Unlike the Agda formalization which often relies on setoids to maintain compatibility with different Agda modes,
Anders is natively cubical. This allows us to use Path types for equality between morphisms, which are
proof-relevant by default. A univalent category in Anders is a precategory where the map from paths between
objects to isomorphisms is an equivalence.

\section{Core Library Structure}

The core library consists of the following modules:
\begin{itemize}
    \item \texttt{category.anders}: Basic definitions of precategories and categories.
    \item \texttt{functor.anders}: Functors and their compositions.
    \item \texttt{natural.anders}: Natural transformations and isomorphisms.
    \item \texttt{adjunction.anders}: Adjunctions between categories.
\end{itemize}

\section{Natural Transformations and Isomorphisms}

Natural transformations are implemented as Sigma types consisting of the component map $\eta$ and a naturality condition.
We have formalized the definitional duality for natural transformations: if $\alpha : F \Rightarrow G$ is a natural
transformation between functors $F, G : \catC \to \catD$, then $\alpha$ also defines a natural transformation
$\alpha\op : G\op \Rightarrow F\op$ between the dual functors. In Anders, this duality is definitional,
meaning $(\alpha\op)\op \equiv \alpha$.

\section{Adjunctions}

Adjunctions $L \dashv R$ between categories $\catC$ and $\catD$ are defined via the unit-counit formulation:
\begin{itemize}
    \item A unit natural transformation $\eta : \text{id}_\catC \Rightarrow R \circ L$.
    \item A counit natural transformation $\varepsilon : L \circ R \Rightarrow \text{id}_\catD$.
    \item Two triangle identities (zig and zag).
\end{itemize}
Due to the current limitations of the Anders compiler regarding named projections within nested Sigma types, the implementation uses numeric indices (e.g., \texttt{u.1} for $\eta$) to ensure stable typechecking.

\section{Equivalences of Categories}

A strong equivalence between two categories $\catC$ and $\catD$ is a pair of functors $F : \catC \to \catD$ and $G : \catD \to \catC$ equipped with natural isomorphisms:
\begin{itemize}
    \item $\eta : \text{id}_\catC \cong G \circ F$
    \item $\varepsilon : F \circ G \cong \text{id}_\catD$
\end{itemize}
We have implemented \texttt{opEquivalence}, which shows that an equivalence $e : \catC \cong \catD$ induces an equivalence $e\op : \catC\op \cong \catD\op$. To ensure definitional duality, we utilize \texttt{symNatIso} to adjust the directions of the unit and counit in the dualized structure.

\section{The Category of Categories}

The category of small precategories, denoted $\mathcal{CAT}$, is formalized as a large category (of type $U_2$). The objects are precategories of type $U_1$, and morphisms are functors of type $U$. While $\mathcal{CAT}$ is naturally a $2$-category, we provide its $1$-categorical structure by specifying identity functors and functor composition. The operation of taking the opposite category is currently provided as a definitional involution on the objects of $\mathcal{CAT}$.

\section{Conclusion}

The Anders categorical library continues to expand by incorporating higher-level structures while maintaining strict adherence to definitional duality and cubical principles. The use of symmetric laws and explicit path-based proofs allows for a clean and performant formalization that avoids the overhead of traditional setoid-based approaches.

\begin{thebibliography}{9}
\bibitem{carette2021}
Jason Z. S. Hu and Jacques Carette.
\textit{Formalizing Category Theory in Agda}.
CPP 2021.
\end{thebibliography}

\end{document}
